
\subsection{Corrections des Exercices}

\begin{correctionbox}[Correction Exercice 1 : Paramètres de $\bar{X}_n$]
$E[\bar{X}_{25}] = \mu = 50$.
$\text{Var}(\bar{X}_{25}) = \frac{\sigma^2}{n} = \frac{100}{25} = 4$.
\end{correctionbox}

\begin{correctionbox}[Correction Exercice 2 : Paramètres de $S_n$]
$E[S_{16}] = n\mu = 16 \times 10 = 160$.
$\text{Var}(S_{16}) = n\sigma^2 = 16 \times 3^2 = 16 \times 9 = 144$.
$\sigma_{S_{16}} = \sqrt{\text{Var}(S_{16})} = \sqrt{144} = 12$.
\end{correctionbox}

\begin{correctionbox}[Correction Exercice 3 : Paramètres de $\bar{X}_n$ (bis)]
$E[\bar{X}_{400}] = \mu = 70$.
$\sigma_{\bar{X}_{400}} = \frac{\sigma}{\sqrt{n}} = \frac{10}{\sqrt{400}} = \frac{10}{20} = 0.5$.
\end{correctionbox}

\begin{correctionbox}[Correction Exercice 4 : Paramètres de $S_n$ (bis)]
$E[S_{64}] = n\mu = 64 \times 0.5 = 32$.
$\text{Var}(S_{64}) = n\sigma^2 = 64 \times 0.01 = 0.64$.
\end{correctionbox}

\begin{correctionbox}[Correction Exercice 5 : Retrouver $\sigma^2$ (via $\bar{X}_n$)]
$\text{Var}(\bar{X}_n) = \frac{\sigma^2}{n} \implies \sigma^2 = n \cdot \text{Var}(\bar{X}_n)$.
$\sigma^2 = 49 \times 2 = 98$.
\end{correctionbox}

\begin{correctionbox}[Correction Exercice 6 : Retrouver $n$ (via $S_n$)]
$\text{Var}(S_n) = n\sigma^2 \implies n = \frac{\text{Var}(S_n)}{\sigma^2}$.
$n = \frac{300}{12} = 25$.
\end{correctionbox}

\begin{correctionbox}[Correction Exercice 7 : CLT pour $\bar{X}_n$ (Queue Droite)]
$\mu=100, \sigma=15, n=36$. $\bar{X}_{36} \approx \mathcal{N}(\mu, \sigma^2/n)$.
$\mu_{\bar{X}} = 100$. $\sigma_{\bar{X}} = \frac{15}{\sqrt{36}} = \frac{15}{6} = 2.5$.
On cherche $P(\bar{X}_{36} > 105)$.
$Z = \frac{105 - 100}{2.5} = \frac{5}{2.5} = 2$.
$P(Z > 2) = 1 - \Phi(2) \approx 1 - 0.9772 = 0.0228$.
\end{correctionbox}

\begin{correctionbox}[Correction Exercice 8 : CLT pour $\bar{X}_n$ (Queue Gauche)]
$\mu=50, \sigma=8, n=64$. $\bar{X}_{64} \approx \mathcal{N}(\mu, \sigma^2/n)$.
$\mu_{\bar{X}} = 50$. $\sigma_{\bar{X}} = \frac{8}{\sqrt{64}} = \frac{8}{8} = 1$.
On cherche $P(\bar{X}_{64} < 48)$.
$Z = \frac{48 - 50}{1} = -2$.
$P(Z < -2) = \Phi(-2) = 1 - \Phi(2) \approx 1 - 0.9772 = 0.0228$.
\end{correctionbox}

\begin{correctionbox}[Correction Exercice 9 : CLT pour $\bar{X}_n$ (Intervalle)]
$\mu=20, \sigma=5, n=100$. $\bar{X}_{100} \approx \mathcal{N}(\mu, \sigma^2/n)$.
$\mu_{\bar{X}} = 20$. $\sigma_{\bar{X}} = \frac{5}{\sqrt{100}} = \frac{5}{10} = 0.5$.
On cherche $P(19 \le \bar{X}_{100} \le 21.25)$.
$Z_1 = \frac{19 - 20}{0.5} = -2$.
$Z_2 = \frac{21.25 - 20}{0.5} = \frac{1.25}{0.5} = 2.5$.
$P(-2 \le Z \le 2.5) = \Phi(2.5) - \Phi(-2) = \Phi(2.5) - (1 - \Phi(2))$.
$P \approx 0.9938 - (1 - 0.9772) = 0.9938 - 0.0228 = 0.9710$.
\end{correctionbox}

\begin{correctionbox}[Correction Exercice 10 : Application $\bar{X}_n$ (Bouteilles)]
$\mu=500, \sigma=6, n=36$. $\bar{X}_{36} \approx \mathcal{N}(\mu, \sigma^2/n)$.
$\mu_{\bar{X}} = 500$. $\sigma_{\bar{X}} = \frac{6}{\sqrt{36}} = \frac{6}{6} = 1$.
On cherche $P(\bar{X}_{36} > 501.5)$.
$Z = \frac{501.5 - 500}{1} = 1.5$.
$P(Z > 1.5) = 1 - \Phi(1.5) \approx 1 - 0.9332 = 0.0668$.
\end{correctionbox}

\begin{correctionbox}[Correction Exercice 11 : Application $\bar{X}_n$ (Tailles)]
$\mu=175, \sigma=8, n=64$. $\bar{X}_{64} \approx \mathcal{N}(\mu, \sigma^2/n)$.
$\mu_{\bar{X}} = 175$. $\sigma_{\bar{X}} = \frac{8}{\sqrt{64}} = \frac{8}{8} = 1$.
On cherche $P(\bar{X}_{64} < 173)$.
$Z = \frac{173 - 175}{1} = -2$.
$P(Z < -2) = \Phi(-2) = 1 - \Phi(2) \approx 1 - 0.9772 = 0.0228$.
\end{correctionbox}

\begin{correctionbox}[Correction Exercice 12 : Application $\bar{X}_n$ (Notes)]
$\mu=70, \sigma=12, n=36$. $\bar{X}_{36} \approx \mathcal{N}(\mu, \sigma^2/n)$.
$\mu_{\bar{X}} = 70$. $\sigma_{\bar{X}} = \frac{12}{\sqrt{36}} = \frac{12}{6} = 2$.
On cherche $P(\bar{X}_{36} < 67)$.
$Z = \frac{67 - 70}{2} = -1.5$.
$P(Z < -1.5) = \Phi(-1.5) = 1 - \Phi(1.5) \approx 1 - 0.9332 = 0.0668$.
\end{correctionbox}

\begin{correctionbox}[Correction Exercice 13 : CLT pour $S_n$ (Queue Droite)]
$\mu=10, \sigma=2, n=100$. $S_{100} \approx \mathcal{N}(n\mu, n\sigma^2)$.
$E[S_{100}] = 100 \times 10 = 1000$.
$\sigma_{S_n} = \sigma\sqrt{n} = 2 \times \sqrt{100} = 2 \times 10 = 20$.
On cherche $P(S_{100} > 1020)$.
$Z = \frac{1020 - 1000}{20} = \frac{20}{20} = 1$.
$P(Z > 1) = 1 - \Phi(1) \approx 1 - 0.8413 = 0.1587$.
\end{correctionbox}

\begin{correctionbox}[Correction Exercice 14 : CLT pour $S_n$ (Queue Gauche)]
$\mu=5, \sigma=4, n=64$. $S_{64} \approx \mathcal{N}(n\mu, n\sigma^2)$.
$E[S_{64}] = 64 \times 5 = 320$.
$\sigma_{S_n} = \sigma\sqrt{n} = 4 \times \sqrt{64} = 4 \times 8 = 32$.
On cherche $P(S_{64} \le 304)$.
$Z = \frac{304 - 320}{32} = \frac{-16}{32} = -0.5$.
$P(Z \le -0.5) = \Phi(-0.5) = 1 - \Phi(0.5)$. (Valeur non fournie).
\end{correctionbox}

\begin{correctionbox}[Correction Exercice 15 : CLT pour $S_n$ (Intervalle)]
$\mu=2, \sigma=3, n=36$. $S_{36} \approx \mathcal{N}(n\mu, n\sigma^2)$.
$E[S_{36}] = 36 \times 2 = 72$.
$\sigma_{S_n} = \sigma\sqrt{n} = 3 \times \sqrt{36} = 3 \times 6 = 18$.
On cherche $P(S_{36} > 90)$.
$Z = \frac{90 - 72}{18} = \frac{18}{18} = 1$.
$P(Z > 1) = 1 - \Phi(1) \approx 1 - 0.8413 = 0.1587$.
\end{correctionbox}

\begin{correctionbox}[Correction Exercice 16 : Application $S_n$ (Ascenseur)]
$\mu=70, \sigma=14, n=49$. $S_{49} \approx \mathcal{N}(n\mu, n\sigma^2)$.
$E[S_{49}] = 49 \times 70 = 3430$.
$\sigma_{S_n} = \sigma\sqrt{n} = 14 \times \sqrt{49} = 14 \times 7 = 98$.
On cherche $P(S_{49} > 3528)$.
$Z = \frac{3528 - 3430}{98} = \frac{98}{98} = 1$.
$P(Z > 1) = 1 - \Phi(1) \approx 1 - 0.8413 = 0.1587$.
\end{correctionbox}

\begin{correctionbox}[Correction Exercice 17 : Application $S_n$ (Rendement)]
$\mu=0.001, \sigma=0.01, n=100$. $S_{100} \approx \mathcal{N}(n\mu, n\sigma^2)$.
$E[S_{100}] = 100 \times 0.001 = 0.1$.
$\sigma_{S_n} = \sigma\sqrt{n} = 0.01 \times \sqrt{100} = 0.01 \times 10 = 0.1$.
On cherche $P(S_{100} > 0.2)$.
$Z = \frac{0.2 - 0.1}{0.1} = \frac{0.1}{0.1} = 1$.
$P(Z > 1) = 1 - \Phi(1) \approx 1 - 0.8413 = 0.1587$.
\end{correctionbox}

\begin{correctionbox}[Correction Exercice 18 : Application $S_n$ (Rendement)]
On utilise les mêmes paramètres : $S_{100} \approx \mathcal{N}(0.1, 0.1^2)$.
On cherche $P(S_{100} < 0)$.
$Z = \frac{0 - 0.1}{0.1} = -1$.
$P(Z < -1) = \Phi(-1) = 1 - \Phi(1) \approx 1 - 0.8413 = 0.1587$.
\end{correctionbox}

\begin{correctionbox}[Correction Exercice 19 : Inverse (Trouver $c$ pour $\bar{X}_n$)]
$\mu=50, \sigma=10, n=100$. $\bar{X}_{100} \approx \mathcal{N}(\mu, \sigma^2/n)$.
$\mu_{\bar{X}} = 50$. $\sigma_{\bar{X}} = \frac{10}{\sqrt{100}} = 1$.
On cherche $c$ tel que $P(\bar{X}_{100} \le c) \approx 0.8413$.
$P(Z \le \frac{c - 50}{1}) = 0.8413$.
D'après la table, le $z$ correspondant est $1.0$.
$\frac{c - 50}{1} = 1 \implies c = 51$.
\end{correctionbox}

\begin{correctionbox}[Correction Exercice 20 : Inverse (Trouver $c$ pour $S_n$)]
$\mu=10, \sigma=3, n=36$. $S_{36} \approx \mathcal{N}(n\mu, n\sigma^2)$.
$E[S_{36}] = 36 \times 10 = 360$.
$\sigma_{S_n} = \sigma\sqrt{n} = 3 \times \sqrt{36} = 18$.
On cherche $c$ tel que $P(S_{36} \le c) \approx 0.0013$.
$P(Z \le \frac{c - 360}{18}) = 0.0013$.
On sait $\Phi(3) \approx 0.9987$, donc $\Phi(-3) = 1 - 0.9987 = 0.0013$.
Le $z$ correspondant est $-3.0$.
$\frac{c - 360}{18} = -3 \implies c = 360 - 3(18) = 360 - 54 = 306$.
\end{correctionbox}

\begin{correctionbox}[Correction Exercice 21 : Inverse (Taille d'échantillon $n$)]
$\mu=0, \sigma=10$. On veut $P(|\bar{X}_n| \le 1) \ge 0.95$.
$P(-1 \le \bar{X}_n \le 1) \ge 0.95$.
Standardisation : $\sigma_{\bar{X}} = \frac{10}{\sqrt{n}}$.
$P\left( \frac{-1 - 0}{10/\sqrt{n}} \le Z \le \frac{1 - 0}{10/\sqrt{n}} \right) \ge 0.95$.
$P\left( \frac{-\sqrt{n}}{10} \le Z \le \frac{\sqrt{n}}{10} \right) \ge 0.95$.
On sait $P(-1.96 \le Z \le 1.96) = 0.95$.
On doit donc avoir $\frac{\sqrt{n}}{10} \ge 1.96$.
$\sqrt{n} \ge 19.6 \implies n \ge (19.6)^2 = 384.16$.
Il faut $n = 385$ au minimum.
\end{correctionbox}

\begin{correctionbox}[Correction Exercice 22 : Inverse (Taille d'échantillon $n$)]
$\mu=100, \sigma=20$. On veut $P(\bar{X}_n \ge 102) \le 0.0228$.
$P(Z \ge z) \le 0.0228$. On sait $P(Z \ge 2) = 1 - \Phi(2) \approx 0.0228$.
Donc on a besoin que notre Z-score soit $\ge 2$.
$Z = \frac{102 - 100}{20 / \sqrt{n}} = \frac{2}{20 / \sqrt{n}} = \frac{2\sqrt{n}}{20} = \frac{\sqrt{n}}{10}$.
On pose $Z \ge 2 \implies \frac{\sqrt{n}}{10} \ge 2$.
$\sqrt{n} \ge 20 \implies n \ge 400$.
Il faut $n = 400$ au minimum.
\end{correctionbox}

\begin{correctionbox}[Correction Exercice 23 : Paramètres (Proportion)]
$X_i \sim \text{Bern}(p)$ avec $p=0.25$.
1.  $\mu = E[X_i] = p = 0.25$.
    $\sigma^2 = \text{Var}(X_i) = p(1-p) = 0.25 \times 0.75 = 0.1875$.
2.  $\hat{p} = \bar{X}_n$. $n=400$.
    $E[\hat{p}] = \mu = 0.25$.
    $\text{Var}(\hat{p}) = \frac{\sigma^2}{n} = \frac{0.1875}{400} \approx 0.00046875$.
\end{correctionbox}

\begin{correctionbox}[Correction Exercice 24 : Calcul (Proportion)]
$p=0.5, n=100$. $\hat{p} \approx \mathcal{N}(p, \frac{p(1-p)}{n})$.
$E[\hat{p}] = 0.5$.
$\text{Var}(\hat{p}) = \frac{0.5 \times 0.5}{100} = \frac{0.25}{100} = 0.0025$.
$\sigma_{\hat{p}} = \sqrt{0.0025} = 0.05$.
On cherche $P(\hat{p} > 0.6)$.
$Z = \frac{0.6 - 0.5}{0.05} = \frac{0.1}{0.05} = 2$.
$P(Z > 2) = 1 - \Phi(2) \approx 1 - 0.9772 = 0.0228$.
\end{correctionbox}

\begin{correctionbox}[Correction Exercice 25 : Marge d'Erreur (Proportion)]
$n=1000, \hat{p}=0.54$.
On estime l'erreur standard $SE = \sqrt{\frac{\hat{p}(1-\hat{p})}{n}}$.
$SE = \sqrt{\frac{0.54 \times (1 - 0.54)}{1000}} = \sqrt{\frac{0.2484}{1000}} \approx 0.01576$.
La marge d'erreur à 95\% est $ME = 1.96 \times SE$.
$ME = 1.96 \times 0.01576 \approx 0.0309$.
(Soit $\pm 3.09\%$).
\end{correctionbox}
